\documentclass[11pt]{article}

\usepackage[margin=1.05in,top=1.2in,bottom=1.15in]{geometry}
\usepackage{specimen}   % shared "Specimen Reclassified" design system, docs/specimen.sty

\setrunninghead{METHOD NOTE / KOSELLECK MACHINE}

\hypersetup{
  pdftitle={Testing the Sattelzeit as a Network Event},
  pdfsubject={A method note on the Koselleck Machine},
}

\begin{document}

\specimencover
  {Testing the Sattelzeit\\as a Network Event}
  {A method note on the Koselleck Machine: measuring whether English
   vocabulary reorganized as a system, and not only word by word, across
   1770-1830.}
  {%
  \coverfield{Repository}{koselleck-networks}
  \coverfield{Document}{docs/method.tex $\rightarrow$ docs/method.pdf}
  \coverfield{Corpus span}{1510-1910, uniform 20-year windows}
  \coverfield{Display resolution}{per-variant, auto-picked (1{,}500-word cap); sweep: 15 settings}
  \coverfield{Compiled}{\today}
  }

\tableofcontents

\newpage

\section{Context}

The historian Reinhart Koselleck argued that the period 1770-1830 (the
\term{Sattelzeit}, German for ``saddle period'') was not simply an era
in which many words happened to change meaning at the same time. He
claimed it was a moment when the whole vocabulary of politics and society
reorganized \emph{as a system}: concepts shifted together, in
relation to one another, rather than drifting independently one at a
time. This has always been argued in prose, from close reading of texts.
It has never been tested computationally.

Ryan Heuser's \emph{Computing Koselleck} took the first step towards such
a test: he trained a separate word-embedding model for each historical
period and measured how far each individual word moves from one period
to the next. He found that this movement is unusually large around
1770-1830, which supports the idea that something real happened in that
window. But his method looks at one word at a time. It cannot say
whether words moved \emph{together}, as a system reorganizing around a
shared structure, or whether each word's shift is its own independent
event that simply happens to cluster in time.

This project adds that second, missing test. Instead of asking how far a
single word moves, it asks whether the groupings among words, the
way the whole vocabulary is organized into clusters of related meaning,
change unusually around the same decades. If Koselleck is right, we
should see the cluster structure itself reorganize at the Sattelzeit, not
just individual words drifting off on their own.

One real example from the built data shows what this looks like
concretely. In the network for 1770-1790, \emph{nation} sits in a
community of civic-administration words (\emph{treasurer, sheriffs,
electors, burgesses}), while \emph{revolution} sits in an unrelated
community of shape and measurement terms (\emph{oval, oblong,
compressed, bushels}). By the very next window, 1790-1810, both words
- along with \emph{constitution} - have moved into a single new
community built around abstract systemic and governmental process
words (\emph{government, arrangements, processes, phenomena}). Three
politically central words, previously scattered across unrelated
communities, land in the same one at exactly the transition into the
French Revolutionary period. This is one concrete instance of the
pattern the \term{migration fraction} (item 6 below) measures in
aggregate across the whole vocabulary - not proof on its own, but a
picture of what ``moved together'' actually looks like in the data.

\section{Method}

The method reuses the same per-period training that Heuser's approach
uses, so the two studies remain directly comparable, and adds a network
layer on top that a word-by-word approach cannot produce.

\begin{enumerate}

  \item \term{Split the corpus into time windows.} The corpus is
    divided into even 20-year windows, currently 1510 to 1910, starting at
    1510 rather than 1500 so that both 1770 and 1830 land exactly on a
    period boundary - the transitions that open and close the Sattelzeit
    are each a literal edge between two periods, not an approximation. The
    windows are uniform across the whole timeline, not just around the
    Sattelzeit: measuring more finely right where reorganization was
    expected and more coarsely everywhere else would bias the method
    toward finding exactly what it was looking for. Heuser's own results
    show a second spike of word-level change around 1905-1925, unrelated
    to the Sattelzeit, a pattern a Sattelzeit-only sampling scheme could
    never have picked up. The span reflects the corpus assembled for this
    project so far; the same pipeline applies unchanged to any other span
    or window size, set entirely in configuration, given a period-dated
    corpus to feed it.

  \item \term{Train a language model per period.} For each time window,
    a model (\term{word2vec}) is trained independently, learning a
    numeric representation of every word from how it is actually used in
    that period's text: words that appear in similar contexts end up with
    similar representations. For example, if ``king'' and ``queen'' both
    tend to appear near words like ``crown'' and ``reign'' in a period's
    text, the model gives them similar numeric representations. Training
    a separate model per period, rather
    than one model updated continuously, is what makes it possible to
    later measure both individual word drift and whole-system
    reorganization: a word's meaning only makes sense relative to the
    other words trained alongside it, in that same period.

  \item \term{Build a word-similarity network.} From those
    representations, a network is built for each period: every word is a
    node, linked to its 15 closest neighbours by \term{cosine
    similarity} (a standard measure of how alike two words'
    representations are). A fixed number of nearest neighbours is used
    instead of a similarity threshold, because at this vocabulary scale a
    fixed threshold produces networks too dense to cluster meaningfully:
    in early testing, on a then-smaller corpus, a 0.3 threshold linked
    roughly one in ten possible word pairs, tens to hundreds of millions
    of edges per period - a density that only gets worse now that
    per-period vocabularies reach up to 170{,}000 words.

    A fixed, hand-curated list of about ninety common function words -
    articles, prepositions, conjunctions, the standard pronoun forms,
    auxiliary verbs, plus early-modern spelling variants (\emph{hath},
    \emph{doth}, \emph{thou}, \emph{ye}, \emph{unto}) that EEBO/ECCO/Evans
    spelling keeps unmodernized - is excluded from the network as nodes,
    along with every single-letter token (always either one of these
    words, a Latin citation abbreviation, or an OCR artifact).

    These words still take part in training, where they help place every
    other word's meaning correctly; they are left out only at this
    network-building stage, since the interest here is in words that
    carry conceptual content. The exclusion is this specific list, not a
    part-of-speech rule - it does not remove every pronoun as a
    grammatical category, only the particular forms named above.

  \item \term{Detect communities.} A \term{community detection}
    algorithm (\term{Leiden}) is run on each network, grouping words
    that are densely connected to one another. In practice, these groups
    function as semantic fields: sets of words that, in that period, are
    used in similar ways to each other more than to the rest of the
    vocabulary.

    This differs from other clustering methods in a basic way. Some
    methods, $k$-means among them, fix a number of groups in advance and
    assign each item to whichever group's average position it sits
    closest to. Leiden does neither: it does not fix the number of groups beforehand,
    and it never compares a word to a group's ``average'', it only ever
    asks a relational question about a word's actual connections. For any
    word, and any group it is already directly connected to, Leiden checks
    whether moving that word into the group would make the network as a
    whole denser \emph{within} groups and sparser \emph{between} them than
    plain chance would predict, given how connected that word and that
    group already are overall. If the move helps, it is made; if not, the
    word stays put or a different move is tried. Repeating this, word by
    word, until no single move can improve things any further, is the
    entire algorithm.

    The number this process is trying to improve at every step has a name,
    \term{modularity} (written $Q$). It compares how connected the words
    placed inside the same group actually are against how connected they
    would be expected to be if the same set of connections had been handed
    out at random instead. A grouping that captures real structure scores
    well above zero; a grouping no better than chance scores around zero;
    a grouping that forces genuinely unrelated words together scores
    \emph{below} zero, actively worse than not grouping at all. Leiden's
    search is nothing more than hill-climbing on this one number.

  \item \term{Repeat detection at several levels of detail.} Community
    detection does not have one single ``correct'' answer: a parameter
    called \term{resolution} controls how many groups appear and how
    large they are. Raising resolution produces more, smaller, more
    tightly-defined groups; lowering it produces fewer, larger, more
    loosely-defined ones - by changing how large a word's own connections
    have to be, relative to chance, before joining a group is judged
    worthwhile. For that reason every network is processed at fifteen fixed
    resolution settings, currently 0.1 to 16.0, and any claim about
    reorganization has to hold up across all of them, not just one
    chosen after the fact.

    \begin{claimbox}
    \claimlabel{Hard rule.} No single setting among the fifteen is ever
    treated as ``the'' result and the rest discarded. The sweep's entire
    purpose is to check that the same conclusion holds broadly across
    settings, not to pick whichever one happens to agree with the
    hypothesis. No finding is reported if it only holds at a single,
    convenient resolution.
    \end{claimbox}

    These fifteen values are calibrated to this project's current
    corpus, not a fixed universal range: the sweep has already been
    extended once before as the corpus grew, raised in two steps to
    reach today's ceiling of 16.0 once the British Library supplement
    pushed per-period vocabularies past what the earlier range could
    meaningfully resolve (that extension's full detail lives in
    ``Choosing the labeling resolution'' below, alongside the related
    display-resolution history). A substantially larger future corpus
    would likely need this range extended again the same way.

    A second, separate resolution decides only which single partition is
    displayed and labeled in the interactive explorer below - a reading
    convenience, with no bearing on which finding gets reported. It is
    picked per period and per region automatically, by a bisection
    search independent of the sweep's fifteen values - starting from
    resolution 1.0 and doubling or halving to bracket, then narrowing,
    the lowest resolution at which no community grows past a
    1{,}500-word cap. Unlike the sweep, it has no fixed ceiling or
    floor: a given period and region can land above 16.0 or below 0.1.
    ``Choosing the labeling resolution'' below covers the mechanism and
    why it replaced an earlier fixed choice.

  \item \term{Compare consecutive periods.} For each pair of
    consecutive periods, the comparison is restricted to words that
    appear in both (each period has its own vocabulary). The best
    possible match between one period's groups and the next period's
    groups is found first, because group numbers do not mean the same
    thing from one period to the next, so the correspondence has to be
    established before anything can be compared.

    Concretely (\code{metrics.py}'s \code{align\_communities}): build a
    contingency table (a grid of counts crossing every old group against
    every new group) counting, for every (old group, new group) pair,
    how many shared words fall in that pair, then solve it as a single
    global assignment problem (the \term{Hungarian algorithm}, a standard
    way to find the best one-to-one matching between two groupings, run
    via SciPy) that picks the one-to-one relabeling of new group numbers
    onto old ones maximizing the total word count matched. This is a
    fixed, deterministic optimum over all groups at once, computed once
    per period pair -- there is no percentage-overlap threshold anywhere
    in it, and no per-group decision made independently of the others.

    A single shared word then counts as \emph{moved} if, under that one
    relabeling, its new group's mapped-back old identity differs from
    the old group it was actually in; otherwise it counts as
    \emph{stayed}. The fraction of shared words that still ended up in a
    different group, despite that best match, is what this project calls
    the \term{migration fraction}. For example, if 100 words appear in
    both period A and period B, and 40 of them end up in a different
    group even after the best possible matching, the migration fraction
    for that pair is 0.40. It is this project's central measurement of
    how much reorganization happened between two periods, computed once
    per period pair at each of the fifteen sweep resolutions -- this is
    the figure item 7 tests the hypothesis against.

    The interactive explorer (Section \ref{sec:webapp}) shows a second,
    separate version of the same measurement, computed the identical
    way but between each period pair's own auto-picked \term{display
    resolution} (item 5) rather than a shared sweep value -- since
    different periods can land on different display resolutions, this
    is its own alignment, not a lookup into the sweep result. This is
    the number, and the per-word ``moved''/``stayed'' arrow, that
    belongs next to a community a reader is actually looking at; it is
    never what a reorganization claim is tested against, which is
    always the sweep-based figure above.

  \item \term{Test whether reorganization concentrates at the
    Sattelzeit.} The hypothesis predicts that the migration fraction
    should be unusually high specifically for the transitions that fall
    inside or close out the 1770-1830 window, compared to every other
    transition across the full 1510-1910 timeline.

    \begin{claimbox}
    \claimlabel{Result.} Averaged across the fifteen-resolution sweep,
    the transitions bracketing the Sattelzeit do rise -- combined
    region: 1750-1770 to 1770-1790 at 0.41, 1770-1790 to 1790-1810 at
    0.45, 1790-1810 to 1810-1830 at 0.52 -- but this is a gradual
    climb, not a spike isolated to the window: 1810-1830 to 1830-1850
    stays elevated at 0.50 rather than falling back afterward, and the
    single highest transition in the entire 1510-1910 series falls
    before the Sattelzeit, not inside it. For a sense of scale, the
    migration fraction across the full 1510-1910 series spans a range
    from transitions where the partition stays roughly consistent to
    ones, like the pre-Sattelzeit peak below, where more than half of
    the shared vocabulary changes group; the Sattelzeit-adjacent numbers
    above sit toward the high end of that range without being its
    maximum. (Exact figures below are the fifteen-resolution average
    for one pipeline run; see Reproducibility under Current limitations
    for how much they can be expected to shift on an independent
    rerun.)

    1690-1710 to 1710-1730 measures 0.55 in the combined corpus and 0.62
    in the British-only split (see the region-split check below) --
    both above every transition the Sattelzeit window itself produces.
    That pre-Sattelzeit peak coincides with a real, isolated drop in
    shared vocabulary between those two periods, consistent with a
    corpus-coverage artifact -- ECCO-TCP thins out precisely in this
    range (see Corpus below) -- rather than a genuine reorganization
    signal, though this has not been confirmed either way.

    Taken together, this measurement does not show a clean, isolated
    reorganization spike at 1770-1830: the pattern is a modest, gradual
    rise that survives the region-split check but does not stand out
    from at least one earlier, likely-artifactual transition. Full
    per-resolution numbers and the complete period series are in
    \code{docs/pipeline\_manual.pdf}'s Stage 8 section.
    \end{claimbox}

  \item \term{Repeat the whole measurement per region.} The corpus is
    not evenly British across the timeline, and the imbalance is worst
    exactly where it matters most (the real numbers are in ``Corpus
    coverage'' under Current limitations, below): a change in the
    cluster structure at that boundary could therefore reflect a
    change in \emph{who was writing} rather than a change in the
    language.

    To separate the two, every step above (training, network
    construction, community detection, the resolution sweep, and the
    migration fraction) is also run on each region's documents alone, in
    addition to the combined corpus. If the reorganization is real, it
    should show up inside the British-only and the American-only
    timelines as well, not only in the mixture of them; if it shows up
    only in the combined run, it is an artifact of corpus composition
    rather than semantic history. Nothing here hardcodes which regions
    exist: a document's region is simply the top-level folder its
    source was filed under, so the same check extends to any other
    split of the corpus.


\end{enumerate}

\subsection{Corpus}

Four sources were considered for the corpus; only the first two are
actually used:

\begin{itemize}
  \item \term{Primary source.} The \term{Text Creation Partnership
    (TCP)}: curated, manually corrected transcriptions with no OCR
    noise, covering 1500-1800 (EEBO-TCP, ECCO-TCP, Evans-TCP).
  \item \term{Supplement.} Since the Sattelzeit extends to 1830 and TCP
    stops at 1800, the 1800-1900 gap is filled by the \term{British
    Library}'s digitised 19th-century books collection (49{,}455 books,
    roughly 65{,}000 volumes, OCR-derived text with catalogue metadata
    already attached, public domain).
  \item \term{Considered and rejected.} Project Gutenberg was considered
    first and dropped: its metadata carries only its own digitisation
    date, not the original print publication date, which makes it
    unusable for a method that buckets every document by the year it
    was actually written.
  \item \term{Considered, not used.} HistWords (pre-trained historical
    embeddings) was identified as a possible independent check against
    this project's own result, but no such comparison has actually
    been run - nothing in the pipeline reads or depends on it.
\end{itemize}

Each source is filed under a region: \code{british} for EEBO-TCP,
ECCO-TCP, and now the British Library supplement; \code{american} for
Evans-TCP, which stops at 1800 along with the rest of TCP. The region
tag is not a claim about two different English-language communities -
for most of this timeline, American print culture and British print
culture share the same language, and treating them as linguistically
distinct would not be defensible.

It tracks archive provenance instead:
EEBO/ECCO and the British Library are one continuous British-print
lineage, Evans is an independently curated American-print archive.
As item 8 above explains, that is what makes the region split a
robustness check rather than a dialect comparison; the combined corpus
stays the headline result either way. A direct consequence of this
design: since the British Library supplement is tagged \code{british},
the \code{american} region has no data past 1800 at all, an honest
asymmetry rather than a gap to paper over.

\subsection{Not a closed dataset}

This project is a method. The corpus is one input to it, fetched
separately and deliberately not vendored: a reader who disagrees with
it, or has a better one, can rerun the whole measurement on their own
text rather than take these numbers on trust. Swapping in a different
corpus is either a folder drop (already-valid TCP P4 XML needs no code
change) or one new reader function beside \code{iter\_xml\_members} in
\code{src/parse\_tcp.py}, emitting the same
\code{\{region, source, doc\_id, year, text\}} record the TCP path
already writes - everything downstream reads only that normalised
record and needs no change. README's Adding a new corpus section gives
the full worked contract, including the exact JSON shape and the two
concrete cases. There is no upload interface, hosted service, or
API-key-driven ingestion: clone the repository, put corpus files in the
right folder, rerun the pipeline scripts in order.

\section{The interactive explorer}
\label{sec:webapp}

The historical argument is complete at this point: every quantitative
claim above, including whether reorganization concentrates at the
Sattelzeit, is produced by \code{src/metrics.py} and needs no labels or
web interface at all. What follows describes optional tooling for
reading and browsing that result, not further evidence for or against
the hypothesis.

Alongside the pipeline, a small web tool, the \term{Koselleck
Machine}, makes the per-period networks directly explorable, without
reading raw numbers. It has four views of the same underlying data.

\runin{Timeline.} The primary view. Type a word and see, in one
continuous strip, which community it belonged to in every period it
appears in. Three encodings carry that information:

\begin{enumerate}[itemsep=0.15em,label=\arabic*.]
  \item the community's plain-English label, with its lane (below) as a
    secondary colour cue;
  \item an explicit gap where the word is simply absent from that
    period's vocabulary;
  \item a marked step wherever the word's community actually changed
    from the period before, versus stayed the same.
\end{enumerate}

Earlier drafts of this tool led with the graph explorer below,
but historian collaborators who reviewed it found a force-directed
diagram hard to read at a glance; what they consistently asked for
instead was exactly this, which domain a word belonged to, over time,
read directly rather than inferred from a picture. Clicking any period
in the strip opens that word's full neighbourhood there, reusing the
word search view below.

\runin{Ask.} A conversational front end to the same measured results,
for a question in plain language instead of a chart: type something
like ``did word meaning change the most around 1770-1830?'' and get an
answer grounded in the same networks, communities, and metrics
described above, never freehand generation. A model answers only by
calling a fixed set of tools that read the built data - it never sees
raw tables, only what those tools return - so a question the data
cannot answer gets an honest ``the structure doesn't show that''
instead of a guess, and every claim is tagged by how reliable that
evidence actually is.

\runin{Word search.} The same underlying data, presented as a plain
table instead of a diagram: type a word and a period, and get back its
closest neighbours by similarity, each neighbour's community, and whether
the searched word's own community changed since the previous period.
Useful when the actual numbers matter more than the picture, and reused
directly by the timeline's own per-period drill-in above.

\runin{Graph explorer.} A period slider drives a force-directed network
diagram: nodes are words, their size reflects how connected they are,
and their color marks which community they belong to in that period. A
search box focuses the view on one word and its neighbourhood; moving the
slider one period at a time shows that neighbourhood evolving, and a
change in a word's color between periods marks a change of community
membership. A side panel explains the method in plain terms and shows,
for the currently selected resolution, how the migration measurement
holds up across all fifteen resolution settings - the same sweep
described above. Kept for browsing the raw network
structure, and because building it first is what surfaced this whole
labeling and lane apparatus, but no longer the primary way to read the
tool, per the reception above.

\runin{Region toggle.} Where the per-region runs described above have been
built, both views also offer a region switch (combined, British only,
American only) so the same word in the same period can be compared across
the sub-corpora directly, and a shift can be checked against the possibility
that it only exists in the mixture. The switch lists only the regions that
were actually built, read off the data itself rather than assumed.

\claimlabel{Interface scope, going forward.} Of the four views above,
Timeline and Ask carry the most weight in practice: Timeline for
tracking one word's story across periods, Ask for querying the
measured findings directly in plain language. The graph explorer gives
comparatively little insight by comparison - a force-directed layout
of hundreds or thousands of nodes doesn't tell a viewer much more than
the same neighbourhood already shown as a table, and it becomes less
useful, not more, exactly as the network grows large enough that the
full graph can no longer be usefully viewed at once. Simplifying the
interface down to Timeline and Ask, dropping or de-emphasizing the
graph explorer, is under consideration.

Every community shown in any view carries a short, plain-English label
(for example ``Government \& Law'', ``Moral Character \& Virtue''), read
by a language model from each community's most representative words: a
convenience for a human reader, not a finding in itself.

A label is not
re-read independently every period. Whenever the same correspondence
step the migration fraction is built on maps a community back to a
predecessor with a label, that label is a candidate to carry forward -
but only after two independent model reads, given the inherited label
and the community's current words, both agree it still fits; either
reader saying no forces a fresh read instead. Across the three current
snapshots (``Where the published labels live'' below), 58--63\% of
labels in the combined and British runs carry forward this
way; in the thinner American run, where more periods are a lineage's own
first appearance, fresh reads are the majority instead.

Two-fifths to a
third of communities in every region are sorted into the
\term{``Structural / Uncertain''} lane rather than a
topical one: the corpus is multilingual (English, Latin, Law French,
Welsh, Scots, and more) and early modern, and a community detected by
this method sometimes groups words by shared grammar (verb
conjugations, comparative forms, spelling variants) rather than by
shared topic. Every view surfaces this lane directly next to the label.

Each community is additionally sorted into one \term{lane}: a small,
fixed list of fifteen broad domains (government and law, religion,
medicine, trade, and so on), plus one explicit ``structural / uncertain''
lane for exactly the ``mixed'' case above. The lane list itself was
authored once by hand, from a read-through of every non-mixed label
already produced, then fixed. A language model assigns each community
to one of these existing lanes; it never invents a new one per run. That
fixed-in-advance property is what lets a lane stay a meaningful,
comparable signal across independent pipeline runs even though, as the
last section explains, the underlying community structure and the raw
label text do not reproduce identically run to run. In the tools above,
a community's plain-English label and number stay the primary signal;
its lane is a secondary colour cue layered on top, not a replacement for
reading the actual label.

The paragraphs above describe what the labels are for. The following
four subsections give the mechanics, so that a technical reader can
see exactly what the model was asked to do and reproduce or replace
it.

\subsection{Choosing the labeling resolution}

Community detection produces a partition at every resolution in the
sweep; the webapp displays and labels only one of them, a choice
separate from the sweep itself. Today that display resolution is
picked automatically: for every period and region-split variant
independently, a bisection search - starting from resolution 1.0 and
doubling or halving to bracket, then narrowing - finds the lowest
resolution at which no community grows past a fixed 1{,}500-word cap.
This search is independent of the fifteen-setting sweep described in
item 5 above, not restricted to picking one of those fifteen values,
so a given period or region can land anywhere, including above 16.0 or
below 0.1. \code{docs/pipeline\_manual.pdf}'s Stage 6 section covers
the search mechanism and the one edge case it took a full production
run to surface.

That automatic rule replaced an earlier, manually-tuned fixed value,
used the same way for every period and region. The display resolution
was originally set to 1.0 project-wide. Adding the British Library
supplement later grew some periods' vocabularies past 170{,}000 words,
and at 1.0 the largest community in the latest periods had grown with
them to 24{,}608 words (14\% of that period's vocabulary) - a size at
which no label, model or human, could describe what was actually
inside it. This was not a new failure mode introduced by the
supplement, only a more visible one: the largest community had already
been a stable 12-15\% of the vocabulary at every period size tested,
small and unremarkable at 70{,}000 words, absurd at 170{,}000.

Modularity falls off smoothly as resolution rises, with no sharp elbow
in the tested range, so there was no purely statistical answer for the
right resolution - the real ceiling was practical, how many communities
per period stay browsable, not mathematical. The fixed value was raised
by a concrete before/after check rather than a formula. Take the
community containing the word ``system'' in 1870-1890, the single worst
offender at the time: at resolution 1.0 its top words read as an
unstructured mix (\emph{perjury, arrogance, indigestion, sensuality,
bronchitis, extortions, dyspepsia}); at 2.0 and 3.0 the same words were
still there, effectively unchanged; only at \textbf{4.0} did the
community actually split, and the half that kept ``system'' became a
real, nameable theme (\emph{arrable, assessments, registration,
payment, loans, forfeitures, coroners, judiciary, treasurers, sheriffs,
freehold, manors, copyhold}) - civil administration and property law.

The same check on two other periods at 4.0 (one small, one mid-sized)
produced equally coherent results (philosophy/rhetoric/mathematics;
church governance and clergy), so the fixed value was raised
project-wide to 4.0 rather than only for the largest periods, and
\code{leiden.resolution\_sweep} was extended to match at the time (up
to 4.0, from 2.0) so the manually-chosen display value and the tested
range stayed aligned.

That manual approach was itself a stopgap: growing the corpus further
would only have called for raising the fixed value again and again,
since the right resolution scales with vocabulary size rather than
being a constant. It was later replaced entirely by the automatic,
per-variant rule described above, and \code{leiden.resolution\_sweep}
was extended again to its current range, fifteen settings from 0.1 to
16.0 (item 5), as part of that same move - so today's display choice
and the tested range stay aligned by construction rather than by manual
bookkeeping.

\subsection{The prompt}

The prompt lives in the repository as a single Markdown file, currently
\pathref{src/prompts/community_labeling_v3.md}, which is parsed at runtime
by \code{src/label\_communities.py} rather than duplicated as a string
constant in code, so the file stays the one place the instructions and
the lane list are edited. It has three sections: the fixed lane list, the
system prompt, and the user message template. One model call is made per
Leiden community, at the auto-picked display resolution described above.

The system prompt, quoted from that file, tells the model:

\begin{quotebox}
``You are labeling word clusters from a computational study of an early
modern through nineteenth-century English-language corpus: the Text
Creation Partnership (EEBO, ECCO, Evans; 1500-1800, clean manually-keyed
transcription) plus the British Library's digitised 19th-century books
collection (1800-1900, OCR text). Each cluster is a `community' found by
running the Leiden algorithm on a per-period word-similarity network,
words that are used in similar contexts more than they are used with the
rest of that period's vocabulary. The corpus is multilingual (English,
Latin, Law French, Welsh, Scots, and other languages appear untranslated)
and pre-modern, so a community sometimes groups words by shared
grammatical form (verb conjugations, comparatives, spelling variants), by
shared foreign language, or by proper names (people, places, biblical
figures) rather than by a real topic. In periods drawn from the British
Library supplement, a community can also be a cluster of OCR fragments -
words with a missing prefix or suffix from a misread hyphen or scan
artifact - rather than real words at all. \emph{All of these are honest
and expected outcomes when they are genuinely what the community is, not
a failure to avoid.}''

It then asks the model to work through, in order, before deciding
anything:

\begin{enumerate}[itemsep=0.3em]
  \item what the community's most central (\emph{Core}) words suggest
    on their own;
  \item whether its less-central (\emph{Mid-rank} and \emph{Peripheral})
    words support or undercut that reading, naming at least one
    specific word from each tier; and
  \item only then, whether a genuine topic holds across all three, or
    whether the community is better described as grammatical,
    foreign-language, or name-driven.
\end{enumerate}

That ordering is itself the point of a required \code{reasoning} field,
a later addition to the prompt made for exactly this reason (see below).
Once that reasoning is done, it asks for exactly three decisions:

\begin{enumerate}[itemsep=0.3em]
  \item ``Whether most of the words shown - across every tier given, not
    just the Core - share a genuine topical or thematic connection a
    historian would recognize (e.g. law, medicine, religion, trade), as
    opposed to being grouped mainly by grammar, a shared non-English
    language, or being mostly proper names or fragments.''
  \item ``A short plain-English label, two to five words, for what the
    community is actually about. If the grouping is grammatical,
    foreign-language, or name-driven rather than topical, the label
    should say what kind of cluster it is (e.g. `Second-Person Verb
    Forms', `Latin Text', `Biblical Names'), not force a topic onto it.''
  \item ``Exactly one lane from the fixed list above. Use `Structural /
    Uncertain' whenever the answer to (1) is no, do not stretch a
    grammatical, foreign-language, or name cluster into a thematic lane
    just because a few of its words are suggestive of one.''
\end{enumerate}
\end{quotebox}

The required \code{reasoning} field is a direct fix for a specific
failure mode, not a general polish pass. An earlier version of this
prompt asked for the same three decisions with no forced reasoning step,
and too often reached for an easy grammatical or ``Structural /
Uncertain'' read from the Core tier alone, without evidence the Mid-rank
and Peripheral words had been weighed at all. Ordering \code{reasoning}
before \code{label} and \code{lane} in the tool schema the model must
call forces that weighing to happen.

A blind check on two real communities, each sent to a fresh,
context-free model twice, once with each version of the prompt, found
the earlier version defaulting to ``Structural / Uncertain'' both times
despite its own draft label already naming a real theme, and the current
version reaching a concrete, evidence-grounded topical lane both times
instead.

The user message carries only the material the model is meant to judge,
filled in from the extraction step: the region, the period, the
community's size in words, and a sample of its member words by
network-degree rank. That sample used to be a flat list of the 25
words with the highest degree in the community, its most central
members, until a real gap surfaced: a flat top-slice only ever shows
words tightly bound to a community's core, hiding a community whose
core looks coherent but whose lower-degree tail has drifted to a
different topic entirely.

The sample is now stratified instead: up to
15 words each from the top (Core), middle (Mid-rank), and bottom
(Peripheral) of the ranked membership list, 45 words in total for any
community large enough to split three ways; a smaller community is
shown in full. The \code{reasoning} field described above exists
specifically so a model has to look at all three tiers, not just the
Core one this older flat sample used to show alone.

Two properties of this setup matter more than the wording. First, the
lane vocabulary is closed. The model is not asked to invent a taxonomy;
it must choose one of fifteen lanes fixed in advance, listed in the same
file and reproduced here in the order the prompt gives them:

\begin{dobox}
\begin{enumerate}[itemsep=0.15em,label=\arabic*.]
  \item Government, Law \& Administration
  \item Religion, Theology \& the Church
  \item Morality: Virtue \& Vice
  \item Medicine, Body \& Health
  \item Science, Mathematics \& Natural Philosophy
  \item Nature, Landscape \& Weather
  \item Military \& Warfare
  \item Trade, Finance \& Commerce
  \item History, Genealogy, Nobility \& Kinship
  \item Rhetoric \& Persuasion
  \item Literature, Drama \& Poetic Diction
  \item Domestic Life, Dress \& Household
  \item Geography \& Territory
  \item Books, Learning \& Scholarship
  \item Structural / Uncertain
\end{enumerate}
\end{dobox}

That list was not invented for the prompt either: it was derived by hand
from 707 labels already produced across the combined, British, and
American runs, then frozen. Second, the response is not free text. The
model must answer by calling a tool, \code{assign\_label}, whose schema
takes a \code{label} string (``two to five word plain-English label'')
and a \code{lane} constrained by a JSON-schema \code{enum} to exactly the
fifteen lanes parsed out of the prompt file. A returned lane outside that
list is rejected and the call retried, so a label with an out-of-taxonomy
lane can never reach the data files.

\subsection{Generating and checking labels}

\code{extract\_community\_words.py} writes, for each populated period at
the configured display resolution, every community's size and its
stratified Core/Mid-rank/Peripheral word sample (see above) to
\pathref{communities/community_words_display[_<region>].json}, once for
the combined corpus and once per region, since a region's own Leiden run
has its own community numbering. \code{label\_communities.py generate}
turns that JSON into a human-editable CSV, walking periods in
chronological order; \code{docs/pipeline\_manual.pdf}'s Stage 7 section
gives the full command sequence and CSV schema, and README's Labeling
communities section has the same commands for a bare clone.

A community whose best correspondence to the period before it (the same
correspondence step the migration fraction is built on) has a label is
a \emph{candidate} to inherit it, not an automatic inheritance: two
independent model reads, given that inherited label and the community's
current words, are asked the same yes-or-no question, does this label
still fit. Only unanimous agreement carries the label forward
unchanged; either reader saying no forces a fresh read, exactly as a
genesis community (no predecessor at all) gets one.

This replaced an
earlier, blunter safeguard, a fixed cap on how many periods a label
could inherit before a forced re-read, itself built after a real case:
a community genuinely holding Dutch and German text at 1510-1530 was
still labeled ``Dutch and German Text'' fifteen periods later, by which
point its content had passed through seven unrelated themes -
including, at the end, English church governance - all unread. Checking
the actual words on every candidate inheritance, rather than counting
periods, catches that drift directly instead of bounding how long it
can go unnoticed.

\code{compile} turns the CSV, including any hand edits, into
\pathref{communities/community_labels_display[_<region>].json}, which is
the exact file \code{webapp/app.py} reads. It also stamps a
\code{\_meta} block into that JSON recording the resolution, the region,
the number of communities, how many were human-edited, how many were
inherited, and a standing caveat that the labels are a single model
read-through and not an empirically validated taxonomy. Compiling used
to append a literal ``(mixed)'' to the label text of every
\code{Structural / Uncertain} community; that is no longer done. The
\code{lane} field already carries the ``this is not a real topic''
signal on its own, and stacking the word ``mixed'' onto an
already-decisive reading (``Second-Person Verb Forms'', ``Biblical Book
Abbreviations'') made even a clear description look like the model
could not settle on one.

\code{publish} copies the CSV and the compiled JSON for a region out of
the data directory and into this repository. \code{{\dd}region} also
accepts a region name or \code{all} at every subcommand.

\subsection{Where the published labels live}

The label files that the pipeline writes live in the data directory,
which is gitignored and does not travel with the repository. That is why
a collaborator cloning the repo previously saw no community names at
all. The \code{publish} subcommand exists to fix exactly that: it copies
a snapshot into the repository's \code{labels/} directory, which
currently holds one CSV and one compiled JSON per region. The filenames
no longer embed a resolution number: since the display resolution is now
picked per variant, not once globally (``Choosing the labeling
resolution'' above), a single number in the filename would be
meaningless once different periods in the same run can land on different
resolutions.

\begin{dobox}
\code{labels/community\_labels\_display.csv}\\
\code{labels/community\_labels\_display.json}\\
\code{labels/community\_labels\_display\_british.csv}\\
\code{labels/community\_labels\_display\_british.json}\\
\code{labels/community\_labels\_display\_american.csv}\\
\code{labels/community\_labels\_display\_american.json}
\end{dobox}

That is the citable snapshot: the version of the labels this document
describes, readable without running the pipeline or holding the corpus.
\code{publish} only copies files to disk; it never commits, so the
snapshot is reviewed by hand before it is added to version control.

\section{Current limitations}

\claimlabel{Corpus coverage.} 1800-1900 is now covered by the British
Library supplement described above, but coverage is uneven in four
ways:

\begin{enumerate}[itemsep=0.15em,label=\arabic*.]
  \item ECCO is thin after 1780, so 1770-1790 leans on more American
    than British documents (Evans-TCP, 1{,}030, against mostly-ECCO's
    792).
  \item The British Library supplement brings the next window,
    1790-1810, close to parity instead.
  \item The \code{american} region has no data at all past 1800, since
    Evans-TCP stops there and the supplement is British-only (see
    Corpus above).
  \item Per-region results carry wider uncertainty than the combined
    headline numbers, and the \code{american} region specifically
    cannot speak to anything past the Sattelzeit's closing edge.
\end{enumerate}

\claimlabel{OCR artifacts.} The British Library text is OCR-derived,
unlike TCP's manually corrected transcription. \code{ocr\_refinement.py}
repairs two classes of damage:

\begin{enumerate}[itemsep=0.15em,label=\arabic*.]
  \item an original line-break hyphen left sitting in the text
    (\emph{``utrum- que''} for \emph{utrumque});
  \item the harder case of the OCR dropping that hyphen entirely and
    leaving a bare space instead (\emph{``par ticulars''} for
    \emph{particulars}), fixed by merging two adjacent words only when
    the merged form is real and at least one half isn't independently
    real on its own - checked by hand against the one period a
    reference wordlist covers, 1790-1810 (142 merges made in the first
    120 changed documents, all confirmed correct).
\end{enumerate}

That wordlist is
built only from that period's own TCP text, so periods entirely after
1800 remain unrepaired and untested by this fix. A different, unrelated
source of noise - French and Welsh text surfacing as its own network
community - is fixed separately, by dropping a document once a language
classifier is at least 90\% confident it isn't English.

\claimlabel{Word-level noise in the network.} Two further patterns pass
through into the trained networks and cluster tightly enough by
co-occurrence to form their own communities, neither yet fixed at the
stage that would need to change (which word becomes a node):

\begin{enumerate}[itemsep=0.15em,label=\arabic*.]
  \item Latin citation-apparatus abbreviations (\emph{capvt},
    \emph{deute}, \emph{regu}), correctly transcribed originals, not
    damage of any kind, but carrying no topical content either.
  \item A distinct, unconfirmed TCP transcription pattern, surfaced
    separately when the labeling resolution was raised, in a handful of
    communities from 1510-1650: a missing ``con-''/``com-'' prefix on
    otherwise-recognizable words (\emph{mytted} for \emph{committed},
    \emph{uersation} for \emph{conversation}, \emph{monwealth} for
    \emph{commonwealth}), most likely a scribal abbreviation mark
    mishandled during text extraction - not confirmed against the raw
    XML, and no fix attempted.
\end{enumerate}

\code{docs/pipeline\_manual.pdf}'s Stage 5 section has more detail on
the first of these.

\claimlabel{Label reliability.} A label is a candidate to inherit from
its predecessor only after two independent model reads agree it still
fits the community's current words (``Generating and checking labels''
above); either
reader saying no forces a fresh read. Across the current label
snapshots, that leaves 37--42\% of labels in the combined and British
runs, and the majority in American, as an independent read for a given
pipeline run. That read is still a single, unvalidated pass by a
language model, not a checked ground truth - useful for orientation, not
for citation.

\claimlabel{Reproducibility.} Rerunning the full pipeline from scratch on
the same corpus does not reproduce an identical result. Community
detection itself is properly seeded and deterministic given a fixed
network; the gap sits one step earlier, in training. Word2Vec training
here uses multiple parallel threads for speed, and that kind of parallel
training is not bit-for-bit repeatable even with a fixed random seed,
because the exact order in which threads update shared values depends on
how the operating system happens to schedule them that run. Small
differences in the trained word vectors cascade into a similarity
network with slightly different edges, which cascades into a
structurally similar but not identical community partition - in one
earlier comparison, two full reruns of the same combined corpus produced
community counts a few apart.

This does not bear on the claim under
test: reorganization is measured by whether \code{migration\_fraction}
concentrates at the Sattelzeit, checked across all fifteen resolutions
within a single run, never by whether a specific community keeps its id
or exact member words across two independently retrained copies of the
pipeline. The specific migration-fraction figures reported in the
Result above (0.41 through 0.62) are exactly this kind of point
estimate: expect them to drift somewhat, not to reproduce bit-for-bit,
on an independent rerun. The one thing meant to stay comparable across independent
reruns is the fixed fifteen-lane list, precisely because it is authored
once by hand and never reinvented per run.

\section{Where this leaves the hypothesis}

Koselleck's claim was that the Sattelzeit was a moment when concepts
reorganized \emph{as a system}, not just individual words drifting
independently. The network test built here finds a real, gradual rise
in migration fraction across the transitions that open and close the
1770-1830 window - a direction consistent with that claim - but not a
spike isolated to the window: the rise continues past 1830 rather than
falling back, and a larger, likely-artifactual rise appears earlier,
at 1690-1730 (item 7). That combination is what ``does not show a
clean, isolated reorganization spike'' means in practice: the evidence
leans toward Koselleck being right about direction, without being
strong enough on its own to call the case closed. The result is
suggestive, not confirming, and rests on this network test alone.

\end{document}
